\documentclass[aspectratio=169]{beamer}
\input{header}
\coursecode{JKSSB}

\title{MS Word Fundamentals}
\subtitle{Lesson 22 --- Interface and Core Document Operations}
\author{JKSSB Computer Awareness}
\date{\today}

\begin{document}
\frame{\titlepage}

\begin{frame}{What MS Word Is}
  \begin{itemize}
    \item \key{MS Word} (Microsoft Word) is a \key{word processor} used to
      create, edit, format, and print text documents.
    \item It is part of the \key{Microsoft Office} suite.
    \item A Word document is saved as a file, commonly with the
      \key{.docx} extension.
    \item The window shows the \key{document area}, the \key{Ribbon}, the
      \key{Quick Access Toolbar}, and the \key{status bar}.
  \end{itemize}
  \begin{block}{The one idea}
    A word processor separates \key{typing} the text from \key{formatting}
    its appearance, so the same words can be shown in many ways.
  \end{block}
\end{frame}

\begin{frame}{The Ribbon and Its Tabs}
  \begin{itemize}
    \item The \key{Ribbon} is the strip of commands across the top of the
      window; it replaced the old menu bars.
    \item Commands are grouped under \key{tabs}: Home, Insert, Design,
      Layout, References, Mailings, Review, and View.
    \item Every tab is divided into \key{groups} (for example, the Home tab
      contains Font, Paragraph, and Styles groups).
    \item Some tabs are \key{contextual}: they appear only when a table or
      picture is selected (Table Tools, Picture Format).
  \end{itemize}
  \begin{alertblock}{Exam point}
    The \key{Home} tab holds the most-used formatting commands: font and
    paragraph tools.
  \end{alertblock}
\end{frame}

\begin{frame}{The Quick Access Toolbar}
  \begin{itemize}
    \item The \key{Quick Access Toolbar (QAT)} is the small row of buttons at
      the top-left corner.
    \item By default it shows \key{Save}, \key{Undo}, and \key{Redo}.
    \item It can be \key{customised} to hold any command you use often.
    \item It can be moved \muted{below the Ribbon} if you prefer.
  \end{itemize}
  \begin{block}{Why it matters}
    The QAT stays visible on every tab, so the commands placed on it are
    always one click away.
  \end{block}
\end{frame}

\begin{frame}{Creating, Opening and Saving Documents}
  \begin{itemize}
    \item The \key{File} tab opens \key{Backstage view}, where New, Open,
      Save, Save As, Print, and Close live.
    \item \key{Ctrl+N} creates a new blank document; \key{Ctrl+O} opens an
      existing one.
    \item \key{Ctrl+S} saves the document; \key{Save As} saves a copy under a
      new name or in a new location.
    \item The default Word format is \key{.docx}; older documents use
      \key{.doc}.
  \end{itemize}
  \begin{alertblock}{Exam point}
    \key{Ctrl+S} = Save. \key{Save As} is used to store a \key{copy} with a
    different name, location, or format.
  \end{alertblock}
\end{frame}

\begin{frame}{Entering and Editing Text}
  \begin{itemize}
    \item New text is typed at the \key{insertion point} (the blinking
      cursor).
    \item \key{Word wrap} moves text to the next line automatically; Enter
      starts a \key{new paragraph}.
    \item \key{Backspace} deletes the character to the left; \key{Delete}
      deletes the character to the right.
    \item \key{Ctrl+Z} undoes the last action; \key{Ctrl+Y} redoes it.
  \end{itemize}
  \begin{block}{Selecting text}
    A \key{word} is selected by double-click; a \key{paragraph} by
    triple-click; the whole document by \key{Ctrl+A}.
  \end{block}
\end{frame}

\begin{frame}{Font Formatting}
  \begin{itemize}
    \item The \key{Font} group on the Home tab changes how characters
      \emph{look}.
    \item It sets the \key{font face} (typeface) and \key{font size}.
    \item Style buttons: \key{Bold (Ctrl+B)}, \key{Italic (Ctrl+I)},
      \key{Underline (Ctrl+U)}.
    \item It also sets \key{font colour}, \key{highlight}, and effects such
      as subscript and superscript.
  \end{itemize}
  \begin{exampleblock}{Worked example}
    To make a heading stand out: select it, choose a larger font size, apply
    \key{Bold}, and pick a darker font colour.
  \end{exampleblock}
\end{frame}

\begin{frame}{Paragraph Alignment}
  \begin{itemize}
    \item \key{Alignment} fixes where each line sits between the margins.
    \item \key{Left} (\key{Ctrl+L}): text lines up on the left --- the
      default.
    \item \key{Centre} (\key{Ctrl+E}): each line is centred.
    \item \key{Right} (\key{Ctrl+R}): text lines up on the right.
    \item \key{Justify} (\key{Ctrl+J}): text is flush at both left and right
      margins.
  \end{itemize}
  \begin{alertblock}{Exam point}
    \key{Justify} stretches spaces so both edges are even; \key{centre} does
    not.
  \end{alertblock}
\end{frame}

\begin{frame}{Indentation}
  \begin{itemize}
    \item \key{Indentation} pushes paragraph lines inward from the margin.
    \item \key{Left indent} and \key{right indent} move the whole paragraph
      inward.
    \item \key{First-line indent}: only the first line is pushed inward.
    \item \key{Hanging indent}: every line \emph{except} the first is pushed
      inward --- used for reference lists.
  \end{itemize}
  \begin{block}{Do not confuse}
    \key{Indentation} shifts lines horizontally; \key{spacing} changes the
    vertical gaps between lines and paragraphs.
  \end{block}
\end{frame}

\begin{frame}{Line and Paragraph Spacing}
  \begin{itemize}
    \item \key{Line spacing} controls the vertical gap between lines within a
      paragraph (for example, single, 1.5, double).
    \item \key{Paragraph spacing} controls the gap \key{before} and
      \key{after} a paragraph.
    \item Both are set from the Home tab's Paragraph group or the Paragraph
      dialog box.
    \item Spacing keeps a document readable without adding blank lines
      manually.
  \end{itemize}
\end{frame}

\begin{frame}{Bullets and Numbering}
  \begin{itemize}
    \item \key{Bullets} mark items with a symbol, for unordered lists.
    \item \key{Numbering} marks items with numbers or letters, for ordered
      lists.
    \item \key{Multilevel lists} combine levels (1, 1.1, 1.1.1) for nested
      outlines.
    \item All three are applied from the \key{Paragraph} group on the Home
      tab.
  \end{itemize}
  \begin{alertblock}{Exam point}
    Use \key{bullets} when order does not matter; use \key{numbering} when
    order or sequence matters.
  \end{alertblock}
\end{frame}

\begin{frame}{Tables}
  \begin{itemize}
    \item A \key{table} arranges text in \key{rows} and \key{columns}; each
      intersection is a \key{cell}.
    \item Insert one from \key{Insert $\rightarrow$ Table}.
    \item Cells can be \key{merged} into one or \key{split} apart.
    \item Rows and columns can be inserted, deleted, and resized.
  \end{itemize}
  \begin{exampleblock}{Worked example}
    A marks table: three columns (Roll No., Name, Marks) and one row per
    student. \key{Table Tools} tabs appear only while the table is selected.
  \end{exampleblock}
\end{frame}

\begin{frame}{Images and Objects}
  \begin{itemize}
    \item \key{Insert $\rightarrow$ Pictures} places an image into the
      document.
    \item Handles on the image let you \key{resize} and rotate it.
    \item \key{Text wrapping} decides how text flows around the image
      (In Line with Text, Square, Tight, Behind Text, and so on).
    \item Charts, shapes, and SmartArt are inserted from the same
      \key{Insert} tab.
  \end{itemize}
  \begin{alertblock}{Exam point}
    ``\key{In Line with Text}'' treats the image like a character; the other
    wrapping styles let text flow around it.
  \end{alertblock}
\end{frame}

\begin{frame}{Headers and Footers}
  \begin{itemize}
    \item A \key{header} is the area at the \key{top} of each page; a
      \key{footer} is the area at the \key{bottom}.
    \item They repeat on \key{every page} of a section.
    \item Insert them from \key{Insert $\rightarrow$ Header / Footer}.
    \item Common uses: document title in the header, \key{page numbers} in
      the footer.
  \end{itemize}
  \begin{block}{Remember}
    Header = top of the page; footer = bottom of the page.
  \end{block}
\end{frame}

\begin{frame}{Page Setup}
  \begin{itemize}
    \item \key{Page setup} is found on the \key{Layout} tab.
    \item \key{Margins}: the blank space between text and the page edges.
    \item \key{Orientation}: \key{Portrait} (tall) or \key{Landscape}
      (wide).
    \item \key{Page size}: for example A4 or Letter.
    \item The same tab sets \key{columns} and page \key{breaks}.
  \end{itemize}
  \begin{alertblock}{Exam point}
    Portrait is taller than wide; Landscape is wider than tall. This is a
    frequent options trap.
  \end{alertblock}
\end{frame}

\begin{frame}{Page Break vs Section Break}
  \begin{itemize}
    \item A \key{Page Break} (\key{Ctrl+Enter}) starts the next content on a
      \key{new page}; formatting remains the same.
    \item A \key{Section Break} starts a new \key{section}, which can have
      its own page setup.
    \item Only a \key{Section Break} lets a part of the document have
      \key{different headers and footers}.
    \item Both are inserted from \key{Layout $\rightarrow$ Breaks}.
  \end{itemize}
  \begin{alertblock}{Trap alert}
    A \key{Page Break} does \key{not} change headers and footers. Changing
    them needs a \key{Section Break}. (TRAP-13)
  \end{alertblock}
\end{frame}

\begin{frame}{Printing a Document}
  \begin{itemize}
    \item \key{Ctrl+P} opens the \key{Print} screen in Backstage view.
    \item The \key{print preview} shows exactly how the pages will appear.
    \item You can choose the \key{printer}, the \key{copies}, and which
      \key{pages} to print.
    \item Clicking \key{Print} sends the document to the selected printer.
  \end{itemize}
  \begin{block}{Before printing}
    Always check the print preview: it catches wrong margins, stray blank
    pages, and cut-off text before paper is used.
  \end{block}
\end{frame}

\begin{frame}{Lesson 22: Recall}
  \begin{itemize}
    \item \key{MS Word} is a word processor; \key{.docx} is its default file
      format.
    \item The \key{Ribbon} carries tabs and groups; the \key{Home} tab holds
      font and paragraph tools.
    \item The \key{Quick Access Toolbar} sits at the top-left and is
      customisable.
    \item \key{Ctrl+S} saves, \key{Ctrl+P} prints, and \key{Ctrl+Enter}
      inserts a page break.
    \item Font tools change character appearance; paragraph tools change
      alignment, indentation, and spacing.
    \item \key{Bullets} are unordered; \key{numbering} is ordered.
    \item A \key{Section Break} (not a Page Break) allows different headers
      and footers. (TRAP-13)
  \end{itemize}
\end{frame}
\end{document}
